\section{Related work}\label{related-work}

\subsection{Reuse, cloning and natural-language
artifacts}\label{reuse-cloning-and-natural-language-artifacts}

Software reuse is studied as component adoption and as the dissemination
of solutions across open ecosystems
\cite{mockus2002two,haefliger2008code}, and large-scale cloning on
GitHub shows that duplicated content is common \cite{lopes2017dejavu};
reuse may occur through dependencies, copying, templates and loosely
standardized processes \cite{heinemann2011extents}. These works address
mainly executable code. In agent skills, behavior is specified in
natural language, and copying does not reveal whether the content was
executed, adapted or adopted independently. We identify identical
contents by hash, as the cloning literature does, but count distinct
\textbf{owners} rather than repositories, to separate diffusion among
developers from duplication by one account. The measure excludes
modified copies; \emph{reuse} is an operational term, implying no
quality or adoption.

\subsection{Agent skills and prompt
artifacts}\label{agent-skills-and-prompt-artifacts}

Research on LLMs with tool use studies selection, planning and execution
of calls
\cite{schick2023toolformer,shen2023hugginggpt,yao2023react,qin2023toolllm}.
An agent skill is instead a persistent artifact, hosted in repositories,
that describes procedures and conventions \cite{agentskills_spec}.
GitSkills provides the corpus \cite{gitskills2027}; other studies
analyze specification and packaging defects
\cite{zhang2026keepskillsreuse}, the effect of skills on real
engineering tasks \cite{han2026sweskillsbench} and the evolution of
prompts \cite{tafreshipour2025prompting}. None asks whether observable
characteristics of a skill and its context are associated with diffusion
across owners, nor how much the construction of the corpus affects that
question.

\subsection{Validation, leakage and construction
bias}\label{validation-leakage-and-construction-bias}

When observations from the same project or organization appear in
training and evaluation, a model may memorize the group instead of
learning relations that generalize \cite{kaufman2012leakage}; the
literature recommends grouped cross-validation
\cite{saeb2017need,roberts2017cross}, and nested validation separates
model selection from error estimation. Models may also exploit
incidental regularities
\cite{geirhos2020shortcut,lapuschkin2019unmasking}. In mined data there
is a variant: attributes produced by the collection or deduplication
procedure itself may depend on the target by construction, without any
column being derived from it. The choice of the representative in
GitSkills is such a case. We treat it as a bias to be audited, and adopt
owner-grouped validation, a per-owner cap, hypotheses fixed before the
holdout and confirmation on disjoint data.

\subsection{Modeling, explainability and research
gap}\label{modeling-explainability-and-research-gap}

We model cross-owner reuse (binary) and \texttt{log1p(owner\_count)}
(RMSLE) with six model families, with and without nested hyperparameter
search. Explainability describes how a model combines attributes but
does not turn associations into causal effects; permutation importance
and partial dependence depend on the model, and correlated variables
share importance
\cite{molnar2020interpretable,shmueli2010explain,breiman2001statistical}.
Our contribution is empirical and methodological for this corpus: to
audit the representative bias, to redefine the target in terms of
owners, and to estimate, with a pre-specified protocol and a new
holdout, how much signal remains.
